\texttt{pfcount} reports, for the process it is given, the number of memory
regions, the total virtual address space, and the minor and major page
faults, and its figures match those the kernel publishes in \texttt{/proc}.
The experiments show the behaviour behind those figures: a mapping costs
virtual space but no memory until it is touched; touching anonymous memory
causes one minor fault per page; reading a file that is not cached causes
major faults; and a cached file is mapped with far fewer faults than pages.

The tool has four limits. The tracepoints are x86-specific, so the
``faults seen'' figures need other hooks on other architectures. The
tracepoints count fault exceptions, not completed faults, so a fault that
ends in \texttt{SIGSEGV} would be seen but not counted as minor or major.
Conversely, a fault the kernel resolves on the process's behalf without an
exception, as when \texttt{MAP\_POPULATE} or \texttt{mlock()} fills a region
in advance, is counted in \texttt{min\_flt} but never seen by the
tracepoints; the programs here do neither, so the two counts agree. Finally, the module looks the tracepoints up by name, which ties it to
kernels that use these names.
